%% ---------------------------------------------------------------------------
%% resumen.tex
%%
%% Resumen (máximo 300 palabras, Lineamientos V4.1) y su traducción fiel al inglés,
%% con el formato del TFG de ejemplo: contexto, qué se desarrolló, cómo, resultados,
%% limitaciones y recomendaciones. Las palabras clave están en \scriptKeywords (main.tex).
%% Cifras: Capítulo 4 (tab:res-modos, tab:res-tiempo-real, tab:res-campo, tab:res-rl).
%%
%% Convención de escritura: una oración por línea del archivo fuente.
%% ---------------------------------------------------------------------------

\chapter*{Resumen}
\thispagestyle{empty}

La localización de un vehículo con sensores de bajo costo enfrenta dos limitaciones: la odometría visual monocular acumula deriva y carece de escala propia, y el GPS de un teléfono tiene errores de varios metros y puede perder la señal.
En este contexto, se desarrolló una plataforma SLAM híbrida que estima en tiempo real la posición de un vehículo combinando la cámara, el GPS y el giroscopio de un teléfono.
La solución incluye una aplicación Android que transmite el video y los sensores, y un canal de procesamiento con odometría visual basada en ORB, escala métrica a partir de profundidad monocular, un filtro de Kalman extendido y una ventana con la trayectoria en vivo.
Además, una política de aprendizaje por refuerzo, entrenada con Q-learning tabular, decide cuánto confiar en la rapidez de la cámara durante las pérdidas del GPS.
El desarrollo siguió una estrategia incremental.
Sobre siete grabaciones en vehículo, con el GPS del teléfono como referencia, la plataforma obtuvo un error absoluto de trayectoria de 2,2~m, frente a 2087~m de la cámara sola.
En una computadora de escritorio procesó 30 cuadros por segundo con una latencia de 107~ms en el percentil 95, y en la prueba de campo procesó en vivo cuatro recorridos en una portátil.
La política, evaluada dejando fuera cada sesión, redujo el percentil 90 del error en los cortes de GPS de 30~s de 170 a 128~m frente al filtro convencional, con un aporte moderado que no se traslada a otro dominio.
Las grabaciones carecen de una referencia de alta precisión, y la escala de la cámara limita la estimación durante las pérdidas largas del GPS.
Se recomienda una referencia más precisa, mejorar la escala a rapideces altas y medir el uso de recursos computacionales.
\bigskip

\textbf{Palabras clave:} \scriptKeywords.

\clearpage
\chapter*{Abstract}
\thispagestyle{empty}

\begin{otherlanguage*}{english}
Locating a vehicle with low-cost sensors faces two limitations: monocular visual odometry accumulates drift and lacks a scale of its own, and a phone's GPS has errors of several meters and can lose the signal.
In this context, a hybrid SLAM platform was developed that estimates the position of a vehicle in real time by combining the camera, the GPS and the gyroscope of a phone.
The solution includes an Android application that streams the video and the sensors, and a processing pipeline with ORB-based visual odometry, metric scale from monocular depth, an extended Kalman filter, and a window with the live trajectory.
In addition, a reinforcement learning policy, trained with tabular Q-learning, decides how much to trust the camera's speed during GPS outages.
The development followed an incremental strategy.
On seven in-vehicle recordings, with the phone's GPS as reference, the platform achieved an absolute trajectory error of 2.2~m, compared with 2087~m for the camera alone.
On a desktop computer it processed 30 frames per second with a latency of 107~ms at the 95th percentile, and in the field test it processed four drives live on a laptop.
The policy, evaluated leaving each session out, reduced the 90th percentile of the error in 30~s GPS outages from 170 to 128~m compared with the conventional filter, with a moderate contribution that does not transfer to another domain.
The recordings lack a high-precision reference, and the camera's scale limits the estimate during long GPS outages.
A more precise reference, a better scale at high speeds, and measuring the use of computational resources are recommended.
\bigskip

\textbf{Keywords:} computer vision, sensor fusion, reinforcement learning, monocular visual SLAM, GPS, real time.
\end{otherlanguage*}

\cleardoublepage

%%% Local Variables:
%%% mode: latex
%%% TeX-master: "main"
%%% End:
